% SOURCE_AUTHORITY: sga5_fr_workpass.tex, lines 12327--12365 (LF-normalized unit).
% SOURCE_SHA256: 791F4EFFC5E02832D5D77ED03518C8156D6F07E4C8238B03545DB93D883FBB28
\subsection*{3. Lembretes sobre as representações lineares dos grupos finitos}

Neste número, $\Lambda$ denota um anel comutativo, e $G$ um grupo.

\subsubsection*{3.1. Caracteres}

Consideremos a subcategoria plena $\mathcal E$ de $\Mod(\Lambda[G])$ dos $\Lambda[G]$-módulos $M$ cujo $\Lambda$-módulo subjacente é projetivo de tipo finito. Para um módulo $M$ deste tipo e para $g\in G$, considera-se o endomorfismo, para a estrutura de $\Lambda$-módulo,
\[
g_M:M\longrightarrow M,
\qquad g_M(x)=gx.
\]
Obtém-se uma função
\[
\chi_M:G\longrightarrow \Lambda,
\qquad \chi_M(g)=\Tr^*_\Lambda(g),
\]
que é uma função central sobre $G$ e chama-se o carácter do $\Lambda[G]$-módulo $M$.

Associando a cada $M\in\Ob(\mathcal E)$ seu caráter $\chi_M$, obtém-se uma função aditiva (VIII 1). Utilizando o sorites dos traços para os endomorfismos de complexos perfeitos, obtém-se até mesmo uma função aditiva $\Tr$ sobre a categoria dos complexos de $\Lambda[G]$-módulos que são perfeitos como complexos de $\Lambda$-módulos. Daí resulta um homomorfismo
\[
\Tr:R^{\bullet}(G)\longrightarrow \mathcal F(G,\Lambda),
\]
onde $R^{\bullet}_{\Lambda}(G)$ é o anel das representações de $G$ sobre $\Lambda$ (VIII 10), e $\mathcal F(G,\Lambda)$ o módulo das funções centrais $G\to\Lambda$.

Quando $G$ é finito, o composto do homomorfismo canônico
\[
K^{\bullet}(\Lambda[G])\longrightarrow R^{\bullet}_{\Lambda}(G)
\]
com $\Tr$ denotar-se-á por $\tr$.

Recorde-se o seguinte resultado (IX th. 1, cor. 3):

\begin{statement}{Proposição 3.2}
Sejam $\mathbb Z_\ell$ o anel dos inteiros $\ell$-ádicos e $G$ um grupo finito. Então o homomorfismo
\[
\tr:K^{\bullet}(\mathbb Z_\ell[G])\longrightarrow \mathcal F(G,\mathbb Z_\ell)
\]
É um monomorfismo. Dois $\mathbb Z_\ell[G]$-módulos projetivos que têm a mesma função traço em $\mathcal F(G,\mathbb Z_\ell)$ são isomorfos.
\end{statement}
